\section{Module Specifications}
\label{sec:module-specifications}

This section specifies the operational workflow of each functional module. Each module is described by its role, processing stages, and interaction contracts.

\subsection{Observer Module}

The Observer is the primary entry point for all requests and handles read-only operations. It follows a ReAct loop: (1)~classify user intent, (2)~select and invoke read-only MCP tools, (3)~chain additional tool calls if evidence is insufficient, (4)~synthesize a grounded response. If a state-changing operation is detected, the Observer produces a structured handoff containing the classified intent, gathered evidence, and target identifiers, then transfers control to the Operator.

The Observer accesses read-only tools spanning queue inspection (\texttt{squeue}), node status (\texttt{sinfo}), job history (\texttt{sacct}), scheduler diagnostics (\texttt{sdiag}, \texttt{sprio}), resource statistics (\texttt{sstat}, \texttt{sreport}), configuration display (\texttt{scontrol\_show}), documentation (\texttt{lookup\_slurm\_docs}), and web search. At runtime, the Observer sees approximately 10 tools from the agent's 18-tool surface; the full MCP server exposes 67 tool endpoints including variants and utilities.

\subsection{Operator Module}

The Operator manages all state-changing operations. It is activated exclusively through Observer handoffs and enforces a Human-in-the-Loop (HITL) confirmation protocol before any destructive action.

\textbf{Workflow}: (1)~Receive handoff with classified intent and targets. (2)~Optionally verify current state with a discovery read. (3)~Present the planned action to the user via a pending action event. (4)~Block until explicit user confirmation or cancellation. (5)~Execute the tool call upon confirmation. (6)~Perform post-execution verification and return results.

\textbf{Tool Classification}: Tools are dynamically classified at startup into three categories---\emph{analysis} (read-only, no confirmation needed), \emph{safe} (information retrieval), and \emph{dangerous} (state-modifying, always requires HITL). Any tool whose name or description implies state modification is classified as dangerous.

\subsection{MCP Tool Layer}

The MCP server provides 67 typed tools organized into functional groups. Table~\ref{tab:mcp-tool-groups} shows the distribution across operational categories. The grouping reflects Slurm's own command families: read-only monitoring tools dominate (suitable for the Observer), while job control, node management, and accounting tools require Operator-level confirmation before execution.

\begin{table}[H]
    \centering
    \caption{MCP Tool Groups}
    \label{tab:mcp-tool-groups}
    \small
    \begin{tabular}{|l|c|l|}
        \hline
        \textbf{Group} & \textbf{Count} & \textbf{Examples} \\
        \hline
        Queue/Monitoring & 10 & \texttt{squeue}, \texttt{sinfo}, \texttt{sacct}, \texttt{sdiag} \\
        Job Control & 8 & \texttt{sbatch}, \texttt{scancel}, \texttt{scontrol\_hold} \\
        Node Management & 6 & \texttt{scontrol\_node}, power up/down \\
        Reservations & 4 & create/delete/update reservation \\
        Accounting & 8 & \texttt{sacctmgr} show/add/modify/delete \\
        Cluster Admin & 5 & reconfigure, shutdown \\
        Documentation \& Web & 3 & \texttt{lookup\_slurm\_docs}, \texttt{web\_search} \\
        Other (triggers, charts, mock) & 23 & triggers, chart generators, test utilities \\
        \hline
        \textbf{Total} & \textbf{67} & \\
        \hline
    \end{tabular}
\end{table}

\textbf{Dual Execution Mode}: The server supports \emph{real mode} (translates tool calls to Slurm CLI commands) and \emph{mock mode} (operates on mutable in-memory state with five predefined scenarios). Mock mode enables deterministic evaluation from known starting states.

\textbf{Parameter Validation}: Before execution, the MCP layer validates job IDs, node names, partition names, and scans batch scripts for dangerous patterns, providing a safety layer independent of LLM reasoning quality.

\subsection{Documentation Retrieval}

A local corpus of 273 official Slurm documentation pages enables grounded answers without relying on LLM parametric memory. The \texttt{lookup\_slurm\_docs} tool performs keyword-based lookup against an indexed corpus. When local docs are insufficient, the agent falls back to \texttt{web\_search} and \texttt{fetch\_web\_content} for external sources.

\subsection{Evaluation Pipeline}

Each benchmark test case specifies: a natural-language request, source cluster state, target state, expected tools, routing decision, and HITL requirement. Scoring uses five weighted dimensions:

\begin{equation}
    \text{Overall} = 0.30 \cdot \text{ToolRecall} + 0.25 \cdot \text{Routing} + 0.25 \cdot \text{HITL} + 0.10 \cdot \text{Keywords} + 0.10 \cdot \text{State}
\end{equation}

A case passes when overall $\geq 0.80$. Aggregate metrics include pass rate, behavioral agreement rate, and safety violation rate. An optional LLM-as-judge provides supplementary 1--5 quality scores without affecting pass/fail determination.
